% appendix: market ecology. Blocks moved intact from Discussion for the
% 9-page NeurIPS main-text boundary (2026-07-19): "Quirk days and the
% crowd", "Who else is in the market", "Negative results". Content
% unchanged except section headers.
\section{Quirk days, the competition, and negative results}
\label{app:market}

\subsection{Quirk days and the crowd}
\label{app:market_quirks}

The true quirk days are a different animal. Roughly once a month the
official print lands opposite to, or within a fraction of a point of, the
direction futures held all night. Futures cannot catch these days at any
hour; that is empirically confirmed, and v1.2 does not pretend to. It
keeps 4 of 12 quirk days safe, right or no-bet at 9:00, exactly as
the baseline does. What it improves is the confidence held on the wrong
side: Mar 10 from 84\% to 66\%, Apr 9 from 84\% to 70\%.

The residual is the auction mechanism itself, the modern descendant of
the stale-open effect \citep{stollwhaley1990}. Even at 9:29 the
futures error floor is $\pm 6$ index points; on Jun 22 (official $-0.14$
points) every model version said $+1.7$ to $+5.7$. Inside the floor, no
price data helps.

The crowd handles these days worst of all. On Mar 10 a bettor posted that
TradingView showed an open of 6{,}792.83 while the official open printed
6{,}796.56, on the other side of the prior close. The retail data feed
and the resolution source disagree exactly when it matters. Hedging with
ETFs or futures fails for the same reason: both track the wrong number,
and Jun 22 showed the ETF leg breaking around a dividend date. Comment
volume tells the rest: median 3 per day, 41 on Apr 9. On these days,
standing down is not a weakness. It is the only correct call futures data
can deliver.

\subsection{Who else is in the market}
\label{app:market_competition}

The competition scan says the edge exists because the informed slots are
empty. The trade tape shows two serious automated actors. One is a print
sniper buying the winner at exactly four seconds after 9:30:00, at
0.989--0.999, in \$1--8k daily clips. The other is a mechanical longshot
harvester that took the cheap winning side on both recent quirk days:
about 1{,}000 trades in three weeks, median buy 15 cents, 93\% of buys
under 50 cents. The 9:25--9:29 buyers at 99 cents and above are
futures confirmers with no informational edge. Nobody shows
auction-informed behavior before the print. On Jun 30 the crowd sat at
35\% on a day that resolved Up, and even the sniper waited.

The insider scan is negative twice over. All 14 wallets with 100\%
day-win rates are stale-book harvesters: average entries at 66--79 cents
on days the model already priced at 0.60--0.72, with hard fills going
about 50\%. The best streak found, 12 of 12 days at a market-implied
joint probability of 2.4\%, is what chance produces among 242 wallets:
about six such streaks are expected. No pre-8:30 release-morning
accumulation appears anywhere, so no BLS-leak signature. The always-right
wallets are manually running our own early-entry stale-book edge at
\$200--8k scale, which validates it and shows the informed slots remain
unoccupied. Against the books directly, the model's Brier score is 0.059
to the crowd's 0.071, worth about 20\% per trade at the 5-cent threshold
over 18 indicative days; on the full 100-day test half the picture is more
even (Table~\ref{tab:baselines}), and the edge lives in the gate and the
quoted price, not in out-forecasting the crowd at every hour.

Hindsight bounds what remains. A perfect oracle over the 58 test days
makes $+\$7{,}315$ against our $+\$164$. More than \$4{,}600 of it comes
from five late quirk-day collapses, and Jun 1 alone, the winning side
at 1.8 cents at 9:29, a 55$\times$, is 82\% of the entire ceiling. A
timing-only oracle on our own 21 days and sides makes $+\$569$, about
3.5$\times$ ours, and it is unharvestable by any tested rule. On 31 of 58
days the best price occurred between midnight and 2:00, exactly when the
model is least confident.

The regime is also moving. The train half's average win multiplier was
1.45 against the test half's 1.28, taker fees of \$0.09--1.02 per trade
arrived on Mar 30, and the profitable entry window has been drifting
earlier. Monthly re-validation is mandatory, not hygiene. Execution
shapes the rest. Our \$20--500 clips on a \$30--48k/day book are the
price impact, so the stale-book edge is structurally a maker strategy
\citep{polymarket2026orders}. A 5-cent modeled edge is not 5 cents of
capturable edge on a taker fill.

\subsection{Negative results}
\label{app:market_negatives}

The catalogue of dead ends is part of the result. On the model side, the
2--3am overnight drift is dead since 2021, so there is no drift term.
Tercile-$k$ without the widened regime sigma is worse than the baseline
it refines, and no single news gather-time works: the decisive moves
spread across the whole night.

The news layer earns its place only as a calendar: about 70\% of ES jumps
occur overnight, and every major US release lands at 8:30. The
deterministic overlay changes almost nothing at current settings (train
\$357 versus \$359, test \$314 versus \$322). The first LLM replay went 2
correct, 1 missed, and 3 wrong-signed out of 6: every wrong sign a
retrieval failure, not a reasoning failure. Retrieval quality dominates
model quality.

On the strategy side, Table~\ref{tab:strategy_graveyard} covers the
families; the designs that died before reaching it are more instructive.
Dip-hunting is refuted by adverse selection. On winning days the median
post-signal dip is 2 cents; on losing days it is 62 cents, the crash to
zero. Every patient variant loses to immediate taker entry on both
halves: dip-triggered entries, resting maker orders one to three cents
inside, 50/50 splits. After a valid signal the market converges
away from you.

Midnight entry is refuted by sizing. Deciding at 00:00 with \$50 stakes
went to ruin on train despite 62\% winners: a 72.5\%-accurate signal
cannot carry a stake. Hour 4 is the earliest hour where confidence
justifies money. Optimal sizing at a small bankroll is all-in by another
name; the compounding rule caps itself at half the balance for exactly
that reason.
